% Figura: Matrix Multiplication Quantum Circuits
% Mostra la struttura riga × colonna completa per un elemento

\begin{figure*}[t]
\centering

% Parte superiore: Visualizzazione riga × colonna
\begin{tikzpicture}[
    cell/.style={draw, minimum size=7mm, anchor=center},
    hlcell/.style={cell, fill=blue!20},
    matlabel/.style={font=\bfseries}
]

% Matrice A
\node[hlcell] (a00) at (0,0) {$a_{00}$};
\node[hlcell, right=-\pgflinewidth of a00] (a01) {$a_{01}$};
\node[cell, below=-\pgflinewidth of a00] (a10) {$a_{10}$};
\node[cell, right=-\pgflinewidth of a10] (a11) {$a_{11}$};
\node[matlabel, above=2mm of a00.north east] {$A$};

% Simbolo ×
\node[right=6mm of a01, font=\Large] (times) {$\times$};

% Matrice B
\node[cell, right=12mm of a01] (b00) {$b_{00}$};
\node[hlcell, right=-\pgflinewidth of b00] (b01) {$b_{01}$};
\node[cell, below=-\pgflinewidth of b00] (b10) {$b_{10}$};
\node[hlcell, right=-\pgflinewidth of b10] (b11) {$b_{11}$};
\node[matlabel, above=2mm of b00.north east] {$B$};

% Simbolo =
\node[right=6mm of b01, font=\Large] (eq) {$=$};

% Matrice C
\node[cell, right=12mm of b01] (c00) {$c_{00}$};
\node[hlcell, right=-\pgflinewidth of c00] (c01) {$c_{01}$};
\node[cell, below=-\pgflinewidth of c00] (c10) {$c_{10}$};
\node[cell, right=-\pgflinewidth of c10] (c11) {$c_{11}$};
\node[matlabel, above=2mm of c00.north east] {$C$};

% Formula esplicita
\node[right=10mm of c01, align=left, font=\small] {
    $c_{01} = a_{00} \cdot b_{01} + a_{01} \cdot b_{11}$\\[2pt]
    \textit{(row 0 $\times$ column 1)}
};

\end{tikzpicture}

\vspace{0.6em}

{\small $\downarrow$ \textit{quantum circuit for computing $c_{01}$}}

\vspace{0.4em}

% Circuito QFT completo
\textbf{(a) QFT-based backend}
\vspace{0.3em}

\begin{tikzpicture}
\begin{yquant*}
    qubit {$\ket{a_{00}}$} a00;
    qubit {$\ket{b_{01}}$} b01;
    qubit {$\ket{a_{01}}$} a01;
    qubit {$\ket{b_{11}}$} b11;
    qubit {$\ket{0}$} acc;

    % MAC #1: a00 × b01 → acc (QFT multiplication + QFT addition)
    box {QFT} acc;
    box {$CC\text{-}R_\phi$} (a00, b01, acc);
    box {$R_\phi$} acc;
    box {IQFT} acc;

    % MAC #2: a01 × b11 → acc
    box {QFT} acc;
    box {$CC\text{-}R_\phi$} (a01, b11, acc);
    box {$R_\phi$} acc;
    box {IQFT} acc;

    output {$\ket{a_{00}}$} a00;
    output {$\ket{b_{01}}$} b01;
    output {$\ket{a_{01}}$} a01;
    output {$\ket{b_{11}}$} b11;
    output {$\ket{c_{01}}$} acc;
\end{yquant*}
\end{tikzpicture}

{\footnotesize Multiplication via controlled-controlled phase rotations ($CC\text{-}R_\phi$); addition via controlled phase rotations ($R_\phi$).}

\vspace{0.5em}

% Circuito Ripple-carry completo
\textbf{(b) Ripple-carry backend}
\vspace{0.3em}

\begin{tikzpicture}
\begin{yquant*}
    qubit {$\ket{a_{00}}$} a00;
    qubit {$\ket{b_{01}}$} b01;
    qubit {$\ket{a_{01}}$} a01;
    qubit {$\ket{b_{11}}$} b11;
    qubit {$\ket{0}$} acc;
    qubit {$\ket{0}$} anc;

    % MAC #1: a00 × b01 → acc (shift-and-add multiplication + ripple addition)
    box {shift} (a00, acc);
    box {C-ADD} (b01, acc, anc);
    box {MAJ} (acc, anc);
    box {$\cdots$} (acc, anc);
    box {UMA} (acc, anc);

    % MAC #2: a01 × b11 → acc
    box {shift} (a01, acc);
    box {C-ADD} (b11, acc, anc);
    box {MAJ} (acc, anc);
    box {$\cdots$} (acc, anc);
    box {UMA} (acc, anc);

    output {$\ket{a_{00}}$} a00;
    output {$\ket{b_{01}}$} b01;
    output {$\ket{a_{01}}$} a01;
    output {$\ket{b_{11}}$} b11;
    output {$\ket{c_{01}}$} acc;
    output {$\ket{0}$} anc;
\end{yquant*}
\end{tikzpicture}

{\footnotesize Multiplication via shift-and-add with controlled ripple addition (C-ADD); addition via majority (MAJ) and unmajority (UMA) gates.}

\vspace{0.4em}
{\small $\vdots$ \textit{repeated for each element $c_{ij}$ using row $i$ of $A$ and column $j$ of $B$}}

\caption{Quantum circuits for matrix multiplication. To compute element $c_{01}$, the circuit performs two multiply-accumulate (MAC) operations using row~0 of $A$ ($a_{00}, a_{01}$) and column~1 of $B$ ($b_{01}, b_{11}$). The two backends differ in both multiplication and addition: (a)~QFT-based uses controlled-controlled phase rotations for multiplication and controlled phase rotations for addition, all operating in the Fourier basis; (b)~Ripple-carry uses shift-and-add for multiplication and Cuccaro's majority/unmajority gates for addition, requiring ancilla qubits. This structure repeats for each output element $c_{ij}$.}
\label{fig:matmul-circuits}
\end{figure*}
